% 量子力学1 量子力学1-16 問題10（135-timeenergy。当時は 量子力学1-9 問題8） 解いた日 2026-11-11
%% kotae
\toi{1}{◎}{}
\[ \delta A\,\delta E=\delta A\,\delta H\ge\frac{\abs{\expval{\comm{\hat{A}}{\hat{H}}}}}{2} \]
\begin{gather*}
\expval{\comm{\hat{A}}{\hat{H}}}=\ii\hbar\frac{\dd\expval{\hat{A}}}{\dd t}\\
\frac{\abs{\expval{\comm{\hat{A}}{\hat{H}}}}}{2}=\frac{\hbar}{2}\abs{\frac{\dd\expval{\hat{A}}}{\dd t}}
\end{gather*}
\[ \therefore\quad\delta A\,\delta E\ge\frac{\hbar}{2}\abs{\frac{\dd\expval{\hat{A}}}{\dd t}} \]
\[ \delta E=\delta H=\sqrt{\expval{\hat{H}^{2}}-\expval{\hat{H}}^{2}} \]
$\hat{H}$ が時間によらないので $\expval{\hat{H}^{2}}$ も $\expval{\hat{H}}$ も時間によらず，$\delta E$ も時間によらない。
\toi{2}{○}{記述}
\[ \tau_{A}=\frac{\delta A}{\abs{\dd\expval{\hat{A}}/\dd t}}\ge\frac{\hbar}{2\delta E} \]
\ayamari{$\hat{A}$ が作用する時間}{$\expval{\hat{A}}$ が $\delta A$ だけ変わる時間}
\naze{$\tau_{A}$ は演算子が作用する時間ではない。$\expval{\hat{A}}$ がその時刻の変化の速さで変わり続けたとき，標準偏差 $\delta A$ だけ変わるのにかかる時間で，$A$ の分布が目に見えて変わる時間の尺度である。}
\[ \therefore\quad\delta E\,\tau_{A}\ge\frac{\hbar}{2} \]
\toi{3}{◎}{}
\[ \hat{H}=\frac{\hat{p}^{2}}{2m} \]
\begin{align*}
\comm{\hat{H}}{\hat{x}}&=\frac{1}{2m}\comm{\hat{p}^{2}}{\hat{x}}\\
&=-\frac{\ii\hbar}{m}\hat{p}
\end{align*}
\begin{align*}
\frac{\dd\expval{\hat{x}}}{\dd t}&=\frac{\ii}{\hbar}\expval{\comm{\hat{H}}{\hat{x}}}\\
&=\frac{\expval{\hat{p}}}{m}
\end{align*}
\[ \tau_{x}=\frac{m\,\delta x}{\abs{\expval{\hat{p}}_{0}}}\qquad\delta x=\sqrt{\expval{\hat{x}^{2}}_{0}-\expval{\hat{x}}_{0}^{2}} \]
\begin{align*}
\expval{\hat{x}^{2}}_{0}&=\int_{-\infty}^{\infty}x^{2}\abs{\psi(x,0)}^{2}\dd x\\
&=(2\pi\sigma^{2})^{-1/2}\int x^{2}\exp\Bigl(-\frac{x^{2}}{2\sigma^{2}}\Bigr)\dd x\\
&=(2\pi\sigma^{2})^{-1/2}\cdot\frac{2\sigma^{2}}{2}\cdot\sqrt{2\pi\sigma^{2}}\\
&=\sigma^{2}
\end{align*}
\begin{align*}
\expval{\hat{x}}_{0}&=\int_{-\infty}^{\infty}x\abs{\psi(x,0)}^{2}\dd x\quad\text{奇関数}\\
&=0
\end{align*}
\[ \expval{\hat{p}}_{0}=\hbar k_{0}\qquad\tau_{x}=\frac{m\sigma}{\hbar k_{0}} \]
\[ \delta E=\sqrt{\expval{\hat{H}^{2}}-\expval{\hat{H}}^{2}} \]
\begin{align*}
\expval{\hat{H}^{2}}&=\frac{1}{4m^{2}}\expval{\hat{p}^{4}}\\
&=\frac{\hbar^{4}}{4m^{2}}\expval{\hat{k}^{4}}
\end{align*}
\begin{align*}
\expval{\hat{H}}&=\frac{1}{2m}\expval{\hat{p}^{2}}\\
&=\frac{\hbar^{2}}{2m}\kakimachigai{\expval{\hat{k}}^{2}}{\expval{\hat{k}^{2}}}
\end{align*}
（脇で $k$ のべきの平均を計算）
\[ \expval{k-k_{0}}=0\qquad\therefore\quad\expval{k}=k_{0} \]
\begin{align*}
&\expval{(k-k_{0})^{2}}\\
&=\expval{k^{2}}-2k_{0}\expval{k}+k_{0}^{2}\\
&=\expval{k^{2}}-k_{0}^{2}\\
&\therefore\quad\expval{k^{2}}=k_{0}^{2}+s^{2}
\end{align*}
\begin{align*}
&\expval{(k-k_{0})^{3}}\\
&=\expval{k^{3}}-3k_{0}\expval{k^{2}}+3k_{0}^{2}\expval{k}-k_{0}^{3}\\
&\therefore\quad\expval{k^{3}}=k_{0}^{3}+3k_{0}(k_{0}^{2}+s^{2})-3k_{0}^{3}\\
&\qquad=k_{0}^{3}+3k_{0}s^{2}
\end{align*}
\begin{align*}
&\expval{(k-k_{0})^{4}}\\
&=\expval{k^{4}}-4k_{0}\expval{k^{3}}+6k_{0}^{2}\expval{k^{2}}\\
&\qquad-4k_{0}^{3}\expval{k}+k_{0}^{4}
\end{align*}
\begin{align*}
\therefore\quad\expval{k^{4}}&=4k_{0}(k_{0}^{3}+3k_{0}s^{2})-6k_{0}^{2}(k_{0}^{2}+s^{2})\\
&\qquad+4k_{0}^{4}-k_{0}^{4}+3s^{4}\\
&=k_{0}^{4}+6k_{0}^{2}s^{2}+3s^{4}
\end{align*}
\[ \delta E=\sqrt{\frac{\hbar^{4}}{4m^{2}}\bigl(\expval{\hat{k}^{4}}-\expval{\hat{k}^{2}}^{2}\bigr)} \]
（脇で引き算）
\begin{align*}
&k_{0}^{4}+6k_{0}^{2}s^{2}+3s^{4}-(k_{0}^{4}+2k_{0}^{2}s^{2}+s^{4})\\
&=4k_{0}^{2}s^{2}+2s^{4}
\end{align*}
\[ \delta E=\frac{\hbar^{2}k_{0}s}{m}\sqrt{1+\frac{s^{2}}{2k_{0}^{2}}} \]
\begin{align*}
\delta E\,\tau_{x}&=\frac{\hbar^{2}k_{0}s}{m}\sqrt{1+\frac{s^{2}}{2k_{0}^{2}}}\cdot\frac{m\sigma}{\hbar k_{0}}\\
&=\hbar\sigma s\sqrt{1+\frac{s^{2}}{2k_{0}^{2}}}\\
&=\frac{\hbar}{2}\sqrt{1+\frac{s^{2}}{2k_{0}^{2}}}\ge\frac{\hbar}{2}
\end{align*}
\toi{4}{◎}{}
\begin{align*}
\expval{E}&=\frac{1}{2}(E_{1}+E_{2})\\
\expval{E^{2}}&=\frac{1}{2}(E_{1}^{2}+E_{2}^{2})
\end{align*}
\begin{align*}
\delta E&=\sqrt{\frac{1}{2}(E_{1}^{2}+E_{2}^{2})-\frac{1}{4}(E_{1}+E_{2})^{2}}\\
&=\frac{1}{2}\sqrt{2E_{1}^{2}+2E_{2}^{2}-E_{1}^{2}-2E_{1}E_{2}-E_{2}^{2}}\\
&=\frac{1}{2}\sqrt{E_{1}^{2}-2E_{1}E_{2}+E_{2}^{2}}\\
&=\frac{1}{2}(E_{2}-E_{1})
\end{align*}
\begin{align*}
\delta E\cdot T&=\frac{1}{2}(E_{2}-E_{1})\cdot\frac{2\pi\hbar}{E_{2}-E_{1}}\\
&=\pi\hbar\ge\frac{\hbar}{2}
\end{align*}
エネルギー固有状態では $\dd\expval{\hat{x}}/\dd t=0$，$\dd\expval{\hat{p}}/\dd t=0$ より
\[ \expval{\hat{p}}=0 \]
\[ \tau_{A}\to\infty\quad\text{発散する} \]
%% kizuki
\item 意味を問われた量は定義を言葉に直す。$\tau_{A}=\delta A/\abs{\dd\expval{\hat{A}}/\dd t}$ は（分布の幅）$\div$（平均の変化の速さ）$=$ 時間で，「$\expval{\hat{A}}$ が $\delta A$ だけ動くのにかかる時間」と読む。
